\documentclass[12pt]{article}
%\usepackage{Sweave}
\usepackage[margin=1in]{geometry}
%\usepackage[parfill]{parskip}
%\usepackage[hidelinks]{hyperref}
\usepackage[colorlinks = true,
            linkcolor = black,
            urlcolor  = blue,
            citecolor = black,
            anchorcolor = black]{hyperref}
\usepackage{rotating}
%\usepackage[backend=biber, style=apa, natbib=true]{biblatex}
\usepackage[citestyle=authoryear, style = apa, natbib=true, backend=bibtex, maxcitenames=2]{biblatex}
\addbibresource{references.bib}
\newtheorem{hyp}{H}
\newtheorem{hypo}{Hypothesis}
\usepackage[compact]{titlesec}
%\usepackage{titlesec}
\titleformat*{\section}{\bfseries\fontsize{14}{18}\selectfont}
\titleformat*{\subsection}{\bfseries\fontsize{12}{16}\selectfont}
%\usepackage{natbib}
\usepackage{csquotes}
\usepackage{xfrac}
\usepackage{mathtools}
\usepackage{booktabs}
\usepackage{lscape}
\usepackage{amsmath}
\usepackage[export]{adjustbox}
\usepackage{graphics}
\usepackage{array}
\usepackage[font=small,skip=1pt]{caption}
\usepackage{subcaption} 
\usepackage{ragged2e}
\usepackage{float}
\usepackage[section]{placeins}
\usepackage{hyperref}
\usepackage{multicol}
\usepackage{setspace}
\usepackage[hang,flushmargin]{footmisc}
\usepackage{afterpage}
\usepackage{soul}
\usepackage{xr}
\usepackage{booktabs}
\usepackage{multirow}
\usepackage{siunitx}
\usepackage{caption}
\usepackage{threeparttable}
\usepackage{makecell} 
\usepackage{tabularx}
\newcommand{\tabitem}{\textbullet~~}
\allowdisplaybreaks

% Set caption to be left-aligned
\captionsetup{justification=raggedright, singlelinecheck=false}

% Configure siunitx
\sisetup{
  group-separator = {,},
  group-minimum-digits = 4 % Start grouping from 4 digits
}

\doublespacing

\title{%
  Invasion in the News: \\
  \large Sentiment, Media Coverage and Signaling}
\author{Donald Moratz and Sarah Gerstein}
\date{\today}

\begin{document}

\maketitle

%\abstract{How do media sources respond to significant disruptions in the international system? This paper examines the reactions of state-controlled media in post-Soviet states and satellite countries to the Russian invasion of Ukraine, comparing them with independent media outlets. Utilizing an extensive corpus of several hundred thousand newspaper articles collected before and after the invasion on February 24, 2022, we employ a large language model trained on adversarially coded data to perform topic classification and sentiment analysis toward Russia. Our findings indicate that sentiment toward Russia becomes significantly more negative in independent media following the invasion, whereas the decline in sentiment is less pronounced in state-owned media. These results highlight the divergent roles of media types in shaping public opinion during international crises and contribute to our understanding of media behavior in times of geopolitical upheaval.}

%\section*{Goal}
%The goal is to write a paper.

%\section*{Research Question}
%How do repeated disasters and their associated magnitudes impact international migration?

\section*{Introduction}

How do media sources respond when there is a significant disruption to the status quo in international system? The Russian invasion of Ukraine in February 2022 was one of the most significant shocks to the international system in the post-Cold War era, the effects of which are still being parsed by researchers and politicians four years later. As the North Atlantic Treaty Organization (NATO) debated how best to prepare for the possibility of an invasion of Ukraine \citep{jakes_for_2024} and a US presidential candidate openly encouraged Russia to invade NATO countries \citep{sanger_outburst_2024}, it is important for researchers to study how states reshaped their relationships with Russia. State media accounts form an important tool for states seeking to influence public opinion \citep{guriev_theory_2020} and evidence suggests that selective attribution and sentiment play an important role in how news is displayed \citep{rozenas_how_2019}. Furthermore, how elite cues impact public opinion is an important topic that has received significant attention within the international relations literature \citep{guisinger_mapping_2017}. On the other hand, the effects of privately owned media on public opinion have received less attention; particularly when state and privately owned media diverge in sentiment. Most previous work has utilized survey experiments to examine the relationship between public opinion and media rather than exploring actual cues from elites. Understanding how strategic elites utilize media messaging to respond to significant events in the international sphere is an important next step in identifying the structure of the relationship between the public and political elites. 

Signalling is another important topic in international relations literature. Much of the theoretical work has centered on how to signal credibility and capability to possible rivals \citep{fearon_signaling_1997}. Other work has suggested that allies may signal by placing troops abroad, prepared and ready to defend alliance commitments \citep{blankenship_trivial_2022}. While the dominant strand of media research within international relations has focused on the role of media in shaping domestic opinions \citep{mccombs_agenda-setting_2002, gershkoff_shaping_2005}, in this article, we argue that the additional role of state-owned media as a signalling device has been overlooked. Building on work by \citet{cohen_theatre_1987}, we argue that states strategically control the sentiment expressed towards Russia in state-owned media to communicate with domestic audiences, but also with other states- signalling either condemnation or a lack thereof of Russia's actions. \citet{cohen_theatre_1987} argues that states, especially USSR aligned states, utilize state-owned media to signal their true intentions when it might be costly to do so through formal means. We argue that the invasion of Ukraine represents a crisis situation where states might be forced by circumstance to say one thing formally, but might hope to signal to other states their true intentions within their media. This strategic behavior stands in stark contrast with independent media which, unconstrained by these strategic considerations, is free to engage in harsh criticism of Russia's actions.

Previous work has suffered from a lack of access to the data necessary to examine these dynamics at scale, but advances in natural language processing offer new avenues for research into the topic. To examine this phenomenon within the real world, we studied how state controlled media in former Soviet states as well as satellite and allied states responded to the Russian invasion of Ukraine compared to independent media sources. To do so, we used an extensive corpus of several hundred thousand newspaper articles collected in the period before and after the invasion on February 24$^{\text{th}}$, 2022. We applied a keyword based approach to identify articles about Russia and then utilized a generative AI model (Chat GPT 4o) to conduct sentiment analysis to examine how the overall sentiment of these articles changes around the discontinuity of the invasion. 

By using a generative AI model, rather than traditional BERT based sentiment classifiers, we are able to isolate the sentiment of the article specifically towards Russia and control for the topic, ensuring that we capture how coverage of Russia changes, separate from potential confounding by a shift in topic towards war-related articles. This approach will allow us to causally identify how strategic elites responded to the crisis situation introduced by the invasion and shed further light on the nature of the relationship between elites and the public as well as the role of public opinion in foreign policy. We present three findings. First, that the Russian invasion of Ukraine caused a significant decline in sentiment in media across all topics, separate from the impact in the shift towards coverage of the war. Second,  we present evidence that there is a differential causal effect of the invasion on independent media versus state owned media, highlighting that state-owned media engaged in strategic management of sentiment towards Russia. Finally, we present evidence that states with security agreements with Russia, despite refusing to publicly condemn the invasion in UN votes or public statements, expressed greater negative sentiment post-invasion than neutral countries. This suggests that those states followed the behavior argued by \citet{cohen_theatre_1987}, saying one thing in formal statements, but using media to communicate their true feelings.


%We find that independent media responds significantly negatively to the Russian invasion of Ukraine. Sentiment towards Russia declines by 0.17 points on a scale of -1 (negative) to 1 (positive), a change of 0.27 standard deviations. Conversely, state-owned media responds significantly less negatively, to the tune of 0.06 points less negative, or 0.10 standard deviations. We argue that the explanation for this difference is strategic. Small states that are both geographically and geo-politcally close to Russia utilize state-owned reporting to influence public opinion but also to signal their lack of opposition to the Russian invasion. In support of this argument, we report that the comparative difference in decline in sentiment is true even in non-military related articles, suggesting that while independent media engages in widespread condemnation of Russia, state-owned media's response is significantly more conservative. 

%However, the lack of condemnation is not uniform. When we disagreggate between states with security agreements prior to the invasion (or in the case of Azerbaijan, signed immediately beforehand) and states that have no security alignment with either Russia or NATO, we find that states with security agreements show a significantly more negative decline in their sentiment towards Russia than neutral states post-invasion. While neutral states show no decline in sentiment towards Russia, states with security agreements show a decline in average sentiment of .10 on the scale of -1 (negative) to 1 (positive), which is a change of roughly .18 standard deviations. This despite the fact that in public statements and UN votes, states with security agreements with Russia routinely declined to condemn the invasion. This provides evidence to support our argument that states strategically utilize their state-owned media to engage in international signalling.

%The article proceeds as follows. We first identify relevant literature on alliance politics and signalling as well as the role of the media. We then discuss the case at hand, with a focus on the Russian invasion of Ukraine as well as which states were used in our analysis. We highlight their relationship to Russia today as part of our review. In our third section, we discuss our theory, identifying hypotheses concerning the role of the media based on foreign policy alignment with Russia, as well as the role of state-owned and privately owned media. In our fourth section, we discuss data, research design and our results. The conclusion assesses next steps as well as the implications of divergent media on public opinion.

\newpage

\printbibliography

%\bibliographystyle{apalike}
%\bibliography{references.bib}

\end{document}